% SEGMENT OLP-0012-S01
% Part: sets-functions-relations
% Chapter: relations
% Section: relations-as-sets

\documentclass[../../../include/open-logic-section]{subfiles}

\begin{document}

\olfileid{sfr}{rel}{set}
\olsection{கணங்களாகத் தொடர்புகள்}

\begin{explain}
\olref[sfr][set][imp]{sec} இல் $\Nat$, $\Int$, $\Rat$, $\Real$ ஆகிய
முக்கியமான கணங்களைக் குறிப்பிட்டோம். இவற்றில் சில கணங்களின்
!!{element}s இடையிலான சுவையான தொடர்புகள் சில உங்களுக்கு நினைவிருக்கும்.
எடுத்துக்காட்டாக, இக்கணங்கள் ஒவ்வொன்றிலும் வழக்கமான ஒரு
\emph{வரிசைத் தொடர்பு} உள்ளது. மேலும், ஒவ்வொரு பொருளும் தனக்குத்தானே
கொண்டிருக்கும், வேறு எந்தப் பொருளுடனும் கொண்டிராத
\emph{ஒன்றே ஆதல்} என்னும் தொடர்பும் உள்ளது.
இன்னும் பல சுவையான தொடர்புகளை நாம் காணவிருக்கிறோம்;
சாத்தியமான தொடர்புகள் அவற்றைவிடவும் அதிகம்.
அவற்றைப் பார்ப்பதற்கு முன், தொடர்புகளை ஒரு சிறப்புவகைக் கணமாகக்
கருதலாம் என்பதைச் சுட்டிக்காட்டித் தொடங்குவோம்.

% SEGMENT OLP-0012-S02
இதற்காக, \olref[sfr][set][pai]{sec} இலிருந்து இரண்டு கருத்துகளை நினைவுகூர்க.
முதலாவது, \emph{வரிசைப்படுத்தப்பட்ட ஜோடி} என்ற கருத்து:
$a$, $b$ தரப்பட்டால் $\tuple{a, b}$ ஐ அமைக்கலாம்.
இங்கு உறுப்புகளின் வரிசை \emph{முக்கியமானது}.
எனவே $a \neq b$ எனில் $\tuple{a, b} \neq \tuple{b, a}$.
(இதனை வரிசைப்படுத்தப்படாத ஜோடிகளுடன், அதாவது $2$-உறுப்புக் கணங்களுடன்,
ஒப்பிடுக; அவற்றில் $\{a, b\}=\{b, a\}$.)
இரண்டாவது, \emph{கார்டீசியன் பெருக்கம்} என்ற கருத்து:
$A$, $B$ கணங்கள் எனில் $A \times B$ ஐ அமைக்கலாம்.
இது $x \in A$, $y \in B$ என்ற நிபந்தனைகளை நிறைவேற்றும்
$\tuple{x, y}$ என்னும் அனைத்து ஜோடிகளின் கணமாகும்.
குறிப்பாக, $A^{2}= A \times A$ என்பது $A$ இலிருந்து பெறப்படும்
அனைத்து வரிசைப்படுத்தப்பட்ட ஜோடிகளின் கணமாகும்.

% SEGMENT OLP-0012-S03
இப்போது ஒரு கணத்தின் மீதான குறிப்பிட்ட தொடர்பைக் கருதுவோம்:
இயல் எண்களின் கணமான $\Nat$ இன் மீதான $<$-தொடர்பே அது.
$n<m$ என்ற நிபந்தனையை நிறைவேற்றும் $\tuple{n, m}$ என்னும் எண்களின்
அனைத்து ஜோடிகளையும் கொண்ட கணத்தைக் கருதுக. அதாவது,
\[
  R=\Setabs{\tuple{n, m}}{n, m \in \Nat \text{ மற்றும் } n<m}.
\]
$n$ என்பது $m$ ஐவிடச் சிறியதாக இருப்பதற்கும்,
$\tuple{n, m}$ என்னும் ஜோடி $R$ இன் உறுப்பாக இருப்பதற்கும்
நெருங்கிய தொடர்பு உள்ளது. அதாவது:
\[
      n<m\text{ அப்போதும் அப்போது மட்டுமே }\tuple{n, m} \in R.
\]
உண்மையில், தகவல் எதையும் இழக்காமல், $R$ என்னும் கணமே
$\Nat$ இன் மீதான $<$-தொடர்பு \emph{எனக்} கருதலாம்.

% SEGMENT OLP-0012-S04
இதேபோல், எண்களுக்கு இடையிலான எந்தத் தொடர்புக்கும்
$\Nat^{2}$ இன் ஓர் உட்கணத்தை அமைக்கலாம்.
மறுதிசையில், எண்களின் ஜோடிகளைக் கொண்ட ஏதாவது ஒரு கணம்
$S \subseteq \Nat^{2}$ தரப்பட்டால், அதற்குரிய தொடர்பு எண்களுக்கு இடையே உள்ளது:
$\tuple{n, m} \in S$ ஆக இருக்கும்போது, அப்போதும் அப்போது மட்டுமே
$n$ என்பது $m$ உடன் கொண்டிருக்கும் தொடர்பே அது.
இது பின்வரும் வரையறையை நியாயப்படுத்துகிறது:
\end{explain}

% SEGMENT OLP-0012-S05
\begin{defn}[இருமத் தொடர்பு]
$A$ என்னும் கணத்தின் மீதான \emph{இருமத் தொடர்பு} என்பது
$A^{2}$ இன் ஓர் உட்கணமாகும்.
$R \subseteq A^{2}$ என்பது $A$ இன் மீதான இருமத் தொடர்பாகவும்
$x, y \in A$ ஆகவும் இருந்தால், $\tuple{x, y} \in R$ என்பதற்குப் பதிலாக
சில சமயங்களில் $Rxy$ (அல்லது $xRy$) என எழுதுகிறோம்.
\end{defn}

% SEGMENT OLP-0012-S06
\begin{ex}
  \ollabel{relations}
இயல் எண்களின் ஜோடிகளைக் கொண்ட $\Nat^{2}$ என்னும் கணத்தை,
பின்வருமாறு இருபரிமாண அணியில் பட்டியலிடலாம்:
\[
  \begin{array}{ccccc}
  \mathbf{\tuple{ 0,0 }} & \tuple{ 0,1 } &
    \tuple{ 0,2 } & \tuple{ 0,3 } & \ldots\\
  \tuple{ 1,0 } & \mathbf{\tuple{ 1,1 }} &
    \tuple{ 1,2 } & \tuple{ 1,3 } & \ldots\\
  \tuple{ 2,0 } & \tuple{ 2,1 } &
    \mathbf{\tuple{ 2,2 }} & \tuple{ 2,3 } & \ldots\\
  \tuple{ 3,0 } & \tuple{ 3,1 } & \tuple{ 3,2 } &
    \mathbf{\tuple{ 3,3 }} & \ldots\\
  \vdots & \vdots & \vdots & \vdots & \mathbf{\ddots}
  \end{array}
\]
இங்கு மூலைவிட்டத்தைத் தடித்த எழுத்தில் காட்டியுள்ளோம்.
ஏனெனில், மூலைவிட்டத்தில் உள்ள ஜோடிகளைக் கொண்ட $\Nat^2$ இன் உட்கணமான
\[
  \{\tuple{0,0 }, \tuple{ 1,1 }, \tuple{ 2,2 }, \dots\},
  \]
என்பது $\Nat$ இன் மீதான \emph{ஒன்றாதல் தொடர்பு} ஆகும்.
(ஒன்றாதல் தொடர்பு அடிக்கடி பயன்படுவதால், எந்த ஒரு கணம் $A$ க்கும்
$\Id{A}=\Setabs{\tuple{ x,x }}{x \in A}$ என வரையறுப்போம்.)
மூலைவிட்டத்திற்கு மேலே உள்ள அனைத்து ஜோடிகளின் உட்கணமான
\[
  L = \{\tuple{ 0,1 },\tuple{ 0,2 },\ldots,\tuple{ 1,2 },
  \tuple{ 1,3 }, \dots, \tuple{ 2,3 }, \tuple{ 2,4 },\ldots\},
\]
என்பது \emph{சிறியது} என்ற தொடர்பாகும்; அதாவது,
$n<m$ ஆக இருக்கும்போது, அப்போதும் அப்போது மட்டுமே $Lnm$.
மூலைவிட்டத்திற்குக் கீழே உள்ள ஜோடிகளின் உட்கணமான
\[
  G=\{ \tuple{ 1,0 },\tuple{ 2,0 },\tuple{
    2,1 }, \tuple{ 3,0 },\tuple{ 3,1 },\tuple{ 3,2 }, \dots\},
\]
என்பது \emph{பெரியது} என்ற தொடர்பாகும்; அதாவது,
$n>m$ ஆக இருக்கும்போது, அப்போதும் அப்போது மட்டுமே $Gnm$.
$L$ உடன் $I$ ஐச் சேர்த்துக் கிடைப்பதை $K=L\cup I$ என அழைக்கலாம்.
இது \emph{சிறியது அல்லது சமமானது} என்ற தொடர்பாகும்:
$n \le m$ ஆக இருக்கும்போது, அப்போதும் அப்போது மட்டுமே $Knm$.
அதேபோல், $H=G \cup I$ என்பது \emph{பெரியது அல்லது சமமானது என்ற தொடர்பு} ஆகும்.
$L$, $G$, $K$, $H$ ஆகிய இத்தொடர்புகள் \emph{வரிசைகள்}
எனப்படும் சிறப்புவகைத் தொடர்புகளாகும்.
எந்த எண்ணும் தனக்குத்தானே $L$ அல்லது $G$ என்ற தொடர்பைக் கொண்டிருப்பதில்லை
என்ற பண்பு $L$, $G$ ஆகியவற்றுக்கு உள்ளது
(அதாவது, அனைத்து $n$ க்கும் $Lnn$, $Gnn$ ஆகிய இரண்டும் பொய்).
இந்தப் பண்பைக் கொண்ட தொடர்புகள் \emph{எங்கும் தற்சுட்டற்றவை} எனப்படும்;
அவை வரிசைகளாகவும் இருந்தால், \emph{கண்டிப்பான வரிசைகள்} எனப்படும்.
\end{ex}

% SEGMENT OLP-0012-S07
\begin{explain}
வரிசைகளும் ஒன்றாதலும் முக்கியமான, இயல்பான தொடர்புகளாக இருந்தாலும்,
நமது வரையறையின்படி $A^{2}$ இன் \emph{எந்த} உட்கணமும்
$A$ இன் மீதான ஒரு தொடர்பாகும் என்பதை வலியுறுத்த வேண்டும்;
அது எவ்வளவு இயல்பற்றதாகவோ செயற்கையானதாகவோ தோன்றினாலும் இதுவே உண்மை.
குறிப்பாக, $\emptyset$ என்பது எந்தக் கணத்தின் மீதும் ஒரு தொடர்பாகும்
(எந்த உறுப்புஜோடியும் கொண்டிராத \emph{வெற்றுத் தொடர்பு} அது).
மேலும், $A^{2}$ என்பதே $A$ இன் மீதான ஒரு தொடர்பாகும்
(ஒவ்வொரு ஜோடியும் கொண்டிருக்கும் தொடர்பு அது);
அது \emph{அனைத்துத் தொடர்பு} எனப்படும்.
ஆனால் $E=\Setabs{\tuple{n, m}}{n>5 \text{ அல்லது } m \times n \ge 34}$
போன்ற ஒன்றும் தொடர்பாகவே கருதப்படுகிறது.
\end{explain}

% SEGMENT OLP-0012-S08
\begin{prob}
$\Pow{\{a, b, c\}}$ என்னும் கணத்தின் மீதான $\subseteq$ தொடர்பின்
!!{element}s[acc] பட்டியலிடுக.
\end{prob}

\end{document}
